% Compilar desde docs/informe con: make pdf
% Fuente autocontenida; no requiere shell-escape ni recursos de red.
\documentclass[11pt,a4paper,oneside]{report}
\usepackage{fontspec}
\setmainfont{Latin Modern Roman}
\setsansfont{Latin Modern Sans}
\setmonofont[Scale=0.88]{Latin Modern Mono Light}
\usepackage{polyglossia}
\setdefaultlanguage{spanish}
\usepackage[margin=2.4cm,headheight=15pt]{geometry}
\usepackage{amsmath,amssymb,booktabs,longtable,tabularx,array,ragged2e}
\usepackage{graphicx,xcolor,tikz}
\usetikzlibrary{arrows.meta,positioning,fit,backgrounds,calc}
\usepackage{enumitem,fancyhdr,titlesec,listings,xurl}
\definecolor{verde}{HTML}{13795B}
\definecolor{tinta}{HTML}{172C26}
\definecolor{gris}{HTML}{60726A}
\definecolor{fondo}{HTML}{F1F5F2}
\usepackage[unicode,colorlinks=true,linkcolor=verde,citecolor=verde,urlcolor=verde]{hyperref}
\hypersetup{pdftitle={Solar Energy Market: proceso de desarrollo y resultados},pdfauthor={Solar Energy Market},pdfsubject={Informe técnico de un prototipo IoT, contabilidad energética y liquidación EVM}}
\newcommand{\revision}{acb2d3192b890e331958c9a6abee155a26747b76}
\newcommand{\corte}{18 de noviembre de 2026}
\newcommand{\code}[1]{{\small\nolinkurl{#1}}}
\newcommand{\repo}[1]{\href{https://github.com/0x3ef/Decentralized-Community-Energy-Network/blob/acb2d3192b890e331958c9a6abee155a26747b76/#1}{\nolinkurl{#1}}}
\newcolumntype{L}[1]{>{\RaggedRight\arraybackslash}p{#1}}
\newcolumntype{Y}{>{\RaggedRight\arraybackslash}X}
\setlength{\parskip}{5pt}
\setlength{\parindent}{0pt}
\sloppy
\setlength{\emergencystretch}{3em}
\renewcommand{\arraystretch}{1.18}
\setlist[itemize]{topsep=3pt,itemsep=2pt,leftmargin=*}
\setlist[enumerate]{topsep=3pt,itemsep=3pt,leftmargin=*}
\titleformat{\chapter}[display]{\sffamily\bfseries\color{tinta}\RaggedRight}{\large\color{verde}Capítulo \thechapter}{6pt}{\Huge}
\titlespacing*{\chapter}{0pt}{-15pt}{22pt}
\titleformat{\section}{\sffamily\Large\bfseries\color{verde}}{\thesection}{.65em}{}
\titleformat{\subsection}{\sffamily\large\bfseries}{\thesubsection}{.65em}{}
\pagestyle{fancy}
\fancyhf{}
\fancyhead[L]{\small\sffamily Solar Energy Market}
\fancyhead[R]{\small\sffamily Proceso y resultados}
\fancyfoot[L]{\footnotesize Corte documental: \corte}
\fancyfoot[R]{\thepage}
\renewcommand{\headrulewidth}{.3pt}
\lstset{basicstyle=\ttfamily\footnotesize,backgroundcolor=\color{fondo},frame=single,rulecolor=\color{fondo},breaklines=true,columns=fullflexible,keepspaces=true,showstringspaces=false,tabsize=2}
\tikzset{box/.style={draw=verde!55,rounded corners=3pt,fill=fondo,align=center,font=\small,minimum height=1cm,text width=3.05cm},arr/.style={-{Latex[length=2mm]},thick,draw=verde},labelnode/.style={font=\scriptsize,color=gris,align=center}}
\begin{document}
\hypersetup{pageanchor=false}
\begin{titlepage}
\color{tinta}
{\sffamily\large COMUNIDAD · ENERGÍA · EVIDENCIA\par}
\vspace{2cm}
{\sffamily\bfseries\Huge Solar Energy Market\par}
\vspace{.65cm}
{\sffamily\LARGE Informe del proceso de desarrollo\\y de los resultados alcanzados\par}
\vspace{1.1cm}
{\large Prototipo de contabilidad y compensación energética comunitaria mediante IoT, SQLModel y contratos inteligentes\par}
\vspace{1.5cm}
\rule{\textwidth}{.8pt}
\vspace{.3cm}
\textbf{Alcance:} implementación y demostración local con SUN y ETH nativo de prueba.\\[.25cm]
\textbf{Fecha de corte:} \corte.\\[.25cm]
\textbf{Revisión técnica:} \code{acb2d31}.\\[.25cm]
\textbf{Repositorio:} \href{https://github.com/0x3ef/Decentralized-Community-Energy-Network}{Decentralized-Community-Energy-Network}.
\vfill
\textbf{Naturaleza del resultado}\par
Se documenta un prototipo funcional en Anvil con mediciones simuladas. No se afirma operación comercial, entrega física individual de electricidad, modificación de facturas, despliegue público verificado ni auditoría de seguridad independiente.
\end{titlepage}
\pagenumbering{roman}
\hypersetup{pageanchor=true}
\chapter*{Resumen ejecutivo}
\addcontentsline{toc}{chapter}{Resumen ejecutivo}
El proyecto desarrolló una plataforma que conecta mediciones de hogares, validación de evidencia energética y liquidación mediante contratos inteligentes. Su propósito es distinguir tres hechos que suelen confundirse: energía observada, derechos aprobados y transferencias económicas confirmadas. El prototipo permite registrar hogares y dispositivos, gestionar comunidades y membresías, calcular intervalos, asignar exportación elegible a demanda y demostrar el intercambio de créditos SUN por ETH de prueba.

La solución utiliza un frontend en React y TypeScript; una API FastAPI con SQLModel, SQLAlchemy y Alembic; PostgreSQL; un broker MQTT Mosquitto; un suscriptor de telemetría; un trabajador de contabilidad y liquidación; Redis y cinco contratos EVM. La fuente normal del demostrador es un simulador coordinado de hogares que publica mensajes identificados como \code{SIMULATED}. Los contratos activos son TimelockController, CommunityRegistry, SUNToken, CommunitySettlement y NativeSettlementVault.

El resultado financiero verificable se conserva en un escenario local: 5 kWh de inyección elegible, 3 kWh asignados, emisión de 5 SUN, entrega y retiro por consumo de 3 SUN, pago y retiro de 0,003 ETH del productor y 0,007 ETH remanentes del depósito del consumidor. La oferta restante es de 2 SUN. Las cantidades describen una ejecución en Anvil; no son ingresos reales ni evidencia de un intercambio físico entre dos viviendas.

El proceso incluyó correcciones de reconexión MQTT, sincronización de límites de intervalos, vinculación de identidades, control de permisos, conciliación de transacciones, persistencia de Anvil y recuperación sin eliminar registros de PostgreSQL. La revisión más reciente corrigió la sobrescritura de credenciales de base de datos en Docker. En esta elaboración se verificaron 44 tablas SQLModel, la cabecera Alembic \code{20261112_local_recovery}, 71 rutas de OpenAPI con 80 operaciones HTTP y 32 pruebas de ciclo de vida aprobadas. Estas comprobaciones no equivalen a una nueva ejecución de todas las suites.

El alcance logrado es una integración local demostrable. Permanecen pendientes la calibración de medidores físicos, el despliegue público verificable, una auditoría externa, el endurecimiento operativo y la integración regulatoria y de facturación. El informe conserva explícitamente esta diferencia entre implementación, evidencia disponible y trabajo futuro.

\textbf{Palabras clave:} comunidad energética, IoT, telemetría, SQLModel, SUN, ERC-20, EVM, gobernanza, custodia, conciliación.
\tableofcontents
\listoffigures
\listoftables
\clearpage
\pagenumbering{arabic}
\section{Presentación del proyecto}
\subsection{Nombre y propósito}
La \textbf{Red Descentralizada de Energía Comunitaria} es una aplicación para registrar energía de hogares y gestionar su asignación y pago dentro de una comunidad. Su código está disponible en el repositorio público \href{https://github.com/0x3ef/Decentralized-Community-Energy-Network}{Decentralized-Community-Energy-Network}. El sistema relaciona mediciones, usuarios, créditos SUN y pagos en ETH \citep{software}.

La validación se realizó con sensores simulados y una red local Anvil. Los resultados descritos corresponden a ese entorno de prueba.

\subsection{Problema abordado}
Un hogar puede generar más energía de la que consume mientras otro necesita energía. Para calcular una compensación se requiere conocer cuánto exportó el productor, cuánto consumió el otro hogar y qué cantidad puede asignarse entre ambos.

El proyecto reúne esos datos y aplica reglas para reconocer la energía de cada participante. El proceso comprende la identificación del medidor, la revisión de lecturas, el cálculo por intervalos, la aprobación del lote y la liquidación del pago. La aplicación permite consultar cada etapa y relacionarla con sus registros \citep{readme,settlement}.

\subsection{Objetivo general}
Desarrollar una plataforma que convierta mediciones de energía de hogares en registros de asignación, créditos SUN y pagos entre participantes mediante contratos inteligentes.

\subsection{Objetivos específicos}
\begin{longtable}{L{.27\linewidth}L{.63\linewidth}}
\caption{Objetivos del proyecto y funciones desarrolladas}\label{tab:objetivos}\\
\toprule Objetivo & Función desarrollada \\ \midrule\endfirsthead
\toprule Objetivo & Función desarrollada \\ \midrule\endhead
Registrar participantes & Gestionar usuarios, wallets, hogares, comunidades, nodos y dispositivos.\\
Recibir mediciones & Publicar datos por MQTT y almacenar lecturas autenticadas.\\
Calcular energía & Obtener diferencias entre contadores y agruparlas en intervalos.\\
Asignar excedentes & Distribuir la exportación elegible entre las demandas de la comunidad.\\
Aprobar lotes & Preparar y ejecutar decisiones de gobernanza con una wallet autorizada.\\
Emitir créditos & Entregar SUN por la inyección aprobada.\\
Liquidar pagos & Reservar SUN y ETH, entregar créditos y acreditar ingresos al productor.\\
Registrar compensación & Quemar los SUN correspondientes al consumo aprobado y entregado.\\
Conservar registros & Mantener lecturas, operaciones y estados para su consulta.\\
\bottomrule
\end{longtable}

\subsection{Participantes}
\textbf{Productor.} Aporta exportación elegible, recibe SUN, los deposita en custodia y cobra ETH cuando se liquida una asignación.

\textbf{Consumidor.} Deposita ETH para financiar la energía asignada, recibe SUN y puede utilizarlos para registrar la compensación de su consumo.

\textbf{Operador.} Consulta la operación de la comunidad, los cálculos y las tareas de procesamiento.

\textbf{Oráculo.} Presenta en blockchain los datos del lote calculado por el servidor. Su cuenta tiene permisos para esa tarea.

\textbf{Gobernanza.} Administra las autorizaciones de la comunidad y aprueba los lotes que habilitan la emisión de créditos \citep{blockchain,security}.

\subsection{Fuentes del informe}
El informe se elaboró a partir de la documentación del repositorio, los modelos de datos, los contratos y los resultados de pruebas conservados. Se revisaron las guías de arquitectura, API, IoT, base de datos, liquidación y despliegue. Las cifras de energía y pagos proceden del escenario local de validación. Los resultados de pruebas se presentan por grupo para conservar el contexto de cada ejecución \citep{architecture,database,iot,verification}.

\chapter{Proceso de desarrollo y decisiones}
\section{Estrategia de implementación}
El historial de cambios muestra una construcción incremental: primero una base de contabilidad y operación, después endurecimiento de IoT y persistencia, posteriormente un protocolo gobernado con ETH nativo y, finalmente, integración de interfaz, pruebas y recuperación. Esta reconstrucción corresponde a cambios observables en Git; no se atribuye una metodología formal Scrum, una duración de sprints ni un equipo de tamaño determinado porque esas evidencias no están documentadas.

\begin{longtable}{L{.18\linewidth}L{.24\linewidth}L{.47\linewidth}}
\caption{Hitos técnicos identificables en el historial.}\label{tab:hitos}\\
\toprule Etapa & Revisiones representativas & Trabajo observado \\ \midrule\endfirsthead
\toprule Etapa & Revisiones representativas & Trabajo observado \\ \midrule\endhead
Base funcional & \code{5c1e686}, \code{9ff670f} & Prototipo comunitario inicial y espacio de trabajo conectado a la API.\\
IoT y semillas & \code{a3add96}, \code{7b4f894}, \code{ba3ea71} & Física coordinada por hogar, reconexión, deduplicación y provisión con modelos SQLModel.\\
Protocolo gobernado & \code{71882e7}, \code{ea79b35} & Registro de comunidades y liquidación nativa con separación de roles y acciones sin firma preparadas por API.\\
Automatización & \code{7c87a83}, \code{9b95d24} & Trabajador de intervalos, operaciones durables, conciliación y procedencia de telemetría v2.\\
Interfaz e integración & \code{99fd0ef}, \code{e061ab0} & Interfaz en español, escenario gobernado y verificación de extremo a extremo local.\\
Documentación & \code{aa44223} & Consolidación de guías y distinción entre protocolo activo y artefactos históricos.\\
Planificación y gobernanza & \code{0fdf5c8}, \code{62b5c7e} & Muestreo alineado con UTC, aprobación manual de lotes y corrección del argumento de intervalo.\\
Operación persistente & \code{2ef7536} & Ciclo Docker unificado, checkpoints atómicos y recuperación sobre la misma base SQL.\\
Escenario desde SQL & \code{46d5da8} & Inserción validada de lecturas y verificación real del ciclo ETH/SUN.\\
Credenciales Docker & \code{acb2d31} & Conservación de credenciales originales y adaptación únicamente del endpoint interno.\\
\bottomrule
\end{longtable}
Los cambios se concentran en registros de Git del 10, 11 y 18 de noviembre de 2026. Las fechas de autor, integración y ejecución no son necesariamente equivalentes; el orden de fases anterior es lógico y técnico, no una estimación de esfuerzo. La revisión base del informe es anterior a la incorporación de este documento.\cite{history}

\section{Decisiones que definieron el resultado}
\subsection{Un único modelo de datos vigente}
Se sustituyó el uso operativo de semillas y diagramas SQL desactualizados por modelos SQLModel y migraciones Alembic. Las semillas crean identidades y dispositivos; no insertan créditos pagados ni transacciones confirmadas ficticias. La provisión gobernada añade transacciones reales de autorización en la red local.

\subsection{Evidencia antes de emisión}
La llegada de una lectura no ejecuta un mint. La emisión requiere intervalos elegibles, un compromiso de lote, revisión, finalización y una prueba válida. Esta decisión evita que un dato arbitrario o una cifra del dashboard se trate como derecho económico confirmado.

\subsection{Custodia y liquidación con fondos reales de la red de prueba}
El protocolo nativo utiliza un vault que reserva tanto los SUN del productor como el ETH del consumidor. Se evita reconocer pagos únicamente incrementando columnas SQL. La transferencia económica ocurre en contrato; el servidor indexa y reconcilia su resultado.

\subsection{Separación de automatización y autoridad}
Los servicios automatizan recepción, cálculo, presentación e indexación. La gobernanza finaliza lotes y los participantes firman depósitos y salidas. En el escenario automatizado local, las firmas de prueba sustituyen las intervenciones manuales, sin convertir esa comodidad en política pública de operación.

\section{Incidencias y correcciones}
\begin{longtable}{L{.25\linewidth}L{.32\linewidth}L{.32\linewidth}}
\caption{Problemas encontrados y respuesta aplicada.}\label{tab:incidencias}\\
\toprule Incidencia & Diagnóstico o efecto & Corrección o tratamiento \\ \midrule\endfirsthead
\toprule Incidencia & Diagnóstico o efecto & Corrección o tratamiento \\ \midrule\endhead
MQTT desconectado, código 4 & Publicación intentada sin sesión aceptada o tras pérdida del broker. & Esperar CONNACK, reconectar con límites, conservar muestra y secuencia, y separar identificadores de clientes.\\
Intervalos sin límites exactos & Iniciar el ciclo en un segundo arbitrario podía desplazar todas las muestras. & Planificación por ranuras UTC y límites contables; no rellenar cortes no observados.\\
Argumento de constructor incorrecto & El simulador no iniciaba al pasar la duración contable al generador equivocado. & Mover la opción a SimulationConfig y verificar el arranque.\\
Red Docker inexistente & Contenedores antiguos referían una red ya reemplazada. & Un único proyecto y archivos Compose; recrear contenedores reteniendo volúmenes.\\
Estado Anvil truncado & El nodo no podía interpretar su archivo de estado. & Checkpoints validados y reemplazo atómico; fallback previo o recuperación explícita.\\
Contratos sin bytecode & La DB retenía direcciones de una cadena perdida; llamadas fallaban. & Cuarentena de proyecciones y reemplazo local auditado, sin replay de pagos.\\
Contraseña PostgreSQL rechazada & Compose forzaba credenciales que podían diferir de las del volumen. & Conservar DATABASE\_URL y adaptar solo host/puerto; no rotar automáticamente el rol.\\
Cuenta MetaMask equivocada & Sesión de aplicación y permisos de la extensión seleccionaban identidades distintas. & Comparar direcciones, verificar propiedad y separar perfiles del navegador.\\
Asignación sin fondos o sin SUN & Saldos en wallet no equivalían a activos depositados. & Exponer estados y exigir depósitos/allowances antes de reservar.\\
\bottomrule
\end{longtable}
Las dos últimas incidencias también se observaron en la demostración manual compartida. Su solución operativa no se presenta como un cambio de código no implementado: la selección automática de cuenta por la extensión sigue requiriendo una experiencia de usuario más clara.

\section{Lecciones del proceso}
La integración requirió tratar la identidad y persistencia de la cadena como parte del dominio. Un chain ID reutilizado no identifica por sí solo la misma red; la génesis y los contratos importan. Asimismo, una prueba aislada de un generador no garantiza un intervalo contable completo, y un endpoint saludable no demuestra que una asignación haya terminado.

El trabajo dejó mecanismos explícitos para retener evidencia fallida en lugar de ocultarla. La consecuencia es que pueden permanecer intervalos rechazados o lotes archivados visibles para auditoría. Esta retención es deliberada: el estado histórico no debe cambiarse manualmente para fabricar una demostración exitosa.

\section{Arquitectura e integración}
\subsection{Organización de componentes}
La solución combina una aplicación modular con procesos de trabajo especializados y contratos EVM. Los contenedores facilitan su despliegue, pero no justifican describirla como un conjunto de microservicios completamente independientes. API, suscriptor y worker comparten modelos y persistencia, y algunas capas consultan directamente mediante sesiones SQLAlchemy \citep{architecture}.

\begin{figure}[htbp]
\caption{Componentes y relaciones principales del sistema}\label{fig:arquitectura}

\centering
\begin{tikzpicture}[node distance=1cm and .55cm]
\node[box] (iot) {Sensores físicos\\o simulador explícito};
\node[box,right=of iot] (mqtt) {Mosquitto\\y suscriptor};
\node[box,right=of mqtt] (db) {PostgreSQL\\SQLModel / Alembic};
\node[box,below=1.5cm of iot] (web) {React y TypeScript\\Wallet del usuario};
\node[box,right=of web] (api) {FastAPI\\Reglas y permisos};
\node[box,right=of api] (worker) {Worker\\Cálculo e indexación};
\node[box,below=1.5cm of api] (evm) {Contratos EVM\\Registro, SUN y vault};
\node[box,left=of evm] (gov) {Gobernanza\\Timelock};
\node[box,right=of evm] (redis) {Redis\\Revocación y actividad};
\draw[arr] (iot)--(mqtt);
\draw[arr] (mqtt)--(db);
\draw[arr] (web)--(api);
\draw[arr] (api)--(db);
\draw[arr] (db)--(worker);
\draw[arr] (worker)--(evm);
\draw[arr] (gov)--(evm);
\draw[arr] (web.south) to[out=-60,in=160] node[labelnode,above,sloped]{firma del participante} (evm.west);
\draw[arr] (api)--(evm);
\draw[arr] (worker)--(redis);
\end{tikzpicture}
\par\noindent\small\textit{Nota.} Elaboración propia a partir del diseño del repositorio. Las flechas simplifican las relaciones y no representan todas las consultas de lectura \citep{architecture}.

\end{figure}

\subsection{Frontend}
React 19 y TypeScript organizan las pantallas; Vite construye los recursos y nginx sirve la aplicación en Docker. TanStack Query maneja consultas y caché; viem participa en las interacciones EVM; Recharts presenta series de energía. Las cantidades monetarias no deben tratarse como flotantes JavaScript: la API devuelve cadenas enteras de wei y el cliente utiliza conversiones y enteros grandes.

El frontend prepara revisión de red, destinatario, valor y argumentos; la extensión EIP-1193 aporta la firma del usuario. Las claves de los participantes permanecen en su wallet. La sesión HTTP y la cuenta conectada deben coincidir con la relación de wallet del servidor para operar fondos propios.

\subsection{API y autorización}
El punto de entrada es \code{src.main:app}, con prefijo \code{/api/v1}. FastAPI valida solicitudes y publica OpenAPI. Los servicios comprueban propiedad de hogares, acceso comunitario y permisos. El usuario se autentica con contraseñas bcrypt y JWT; la wallet se verifica con un reto EIP-191 de un solo uso.

La API prepara datos de la llamada al contrato sin firmar. La preparación devuelve los datos que el usuario revisa antes de firmar. Los hashes posteriormente recibidos se siguen y contrastan con la cadena. Las rutas y operaciones de la API se resumen en el apartado de resultados \citep{api}.

\subsection{Recepción y trabajo asíncrono}
El suscriptor ejecuta el bucle MQTT y coordina una cola acotada con el procesamiento asíncrono SQL. Confirma mensajes después de la transacción de persistencia. La deduplicación hace que una redelivery por QoS no sume energía dos veces.

El trabajador procesa contadores, finaliza intervalos, descubre lotes, gestiona outbox, transmite operaciones y audita confirmaciones. Redis aporta latidos y revocación de JWT; no es el libro mayor de pagos. El estado de las tareas se conserva para seguimiento y revisión explícita.

\subsection{Registro de transacciones e indexación}
La modificación de negocio y su evento outbox se confirman conjuntamente. Antes de transmitir una operación del oráculo se conservan sus bytes firmados, hash, nonce y origen. Un bloqueo asesor SQL coordina nonces por chain ID, génesis y remitente. Reintentar una respuesta ambigua no debe crear otro pago con un nonce nuevo.

Una operación se refleja como confirmada después de comprobar recibo exitoso, bloque canónico, profundidad y eventos esperados. Las reorganizaciones y los eventos huérfanos exigen reconstrucción de proyecciones o revisión; la aplicación no inventa una reversión económica de lo ya pagado \citep{architecture,blockchain}.

\subsection{Responsabilidades de los componentes}
Los dispositivos aportan las lecturas. El servidor comprueba su origen y calcula la energía. La gobernanza aprueba el lote y los contratos verifican las pruebas y las condiciones económicas. PostgreSQL conserva la evidencia y la aplicación consulta los resultados confirmados de la cadena.

\section{Infraestructura IoT y contabilidad energética}
\subsection{Modelo del hogar simulado}
La simulación utiliza un estado coordinado por hogar en \code{iot/simulator/physics.py}. Los sensores no inventan curvas independientes que puedan contradecirse: generación, carga, almacenamiento e intercambio se derivan de una misma evolución física sintética. Existen perfiles de consumidor, productor solar, prosumidor, prosumidor con batería y pequeña instalación comercial. Las capacidades instaladas determinan qué canales se publican \citep{iot}.

El balance de potencia del bus de corriente alterna se expresa como:
\begin{equation}
P_{PV}+P_{\text{descarga}}+P_{\text{importación}}=P_{\text{carga}}+P_{\text{batería}}+P_{\text{exportación}}.
\label{eq:balance}
\end{equation}
Las eficiencias de conversión fotovoltaica y almacenamiento se modelan por separado. La batería respeta límites de estado de carga. La integración utiliza tiempo UTC transcurrido y subpasos físicos acotados; no se confunden segundos de espera del programa con tiempo simulado.

La elevación solar depende de fecha y ubicación de Cali; las nubes y las cargas varían gradualmente. Noches sin exportación pueden ser un comportamiento esperado. Estos perfiles permiten pruebas reproducibles, pero no constituyen predicciones meteorológicas ni datos calibrados de una instalación.

\begin{longtable}{L{.22\linewidth}L{.35\linewidth}L{.32\linewidth}}
\caption{Sensores y utilidad de sus datos.}\label{tab:sensores}\\
\toprule Sensor & Datos principales & Uso en la plataforma \\ \midrule\endfirsthead
\toprule Sensor & Datos principales & Uso en la plataforma \\ \midrule\endhead
Medidor inteligente & Contadores de consumo e importación/exportación; potencia, corriente, voltaje y frecuencia. & Fuente principal de intercambio de red y de evidencia para intervalos.\\
Inversor solar & Potencia y energía fotovoltaica, condiciones y temperatura del inversor. & Seguimiento de generación y coherencia con la operación del hogar; no sustituye al medidor de exportación.\\
Monitor de batería & Estado de carga, capacidad, carga/descarga y contadores. & Seguimiento del almacenamiento y balance energético.\\
Estación meteorológica & Irradiancia y temperaturas del entorno y panel. & Contexto de producción solar y diagnóstico de la simulación.\\
\bottomrule
\end{longtable}

\subsection{Modos de simulación y procedencia}
El modo normal es \code{REALTIME}, velocidad uno y fallos desactivados. La publicación registrada se realiza mediante MQTT v2. Los modos deterministas o acelerados permiten auditoría y pruebas acotadas, pero no aportan evidencia de liquidación del flujo normal. El gateway gRPC v1 permanece como ruta diagnóstica, no como transporte principal del simulador gobernado.

El protocolo de mensajes declara \code{source=SIMULATED}; la fuente autorizada del dispositivo en el servidor prevalece sobre lo que afirme el publicador. Cambiar una etiqueta del JSON no convierte una medición virtual en física. Los modos de ensayo permiten evaluar continuidad de publicación, consistencia temporal, identidad de las fuentes y comportamiento de los contadores bajo condiciones controladas.

\subsection{Sobre de telemetría y autenticación}
El tópico canónico es \code{solar/<community>/<node>/<device>/telemetry}. El sobre v2 incluye versión, identificador de mensaje, secuencia, hora de medición, hogar, dispositivo público, comunidad, nodo, fuente, modo, unidades y datos. La credencial se verifica, pero se excluye del payload persistido y de los digests de auditoría.

La aceptación exige correspondencia entre tópico, identidades, dispositivo activo, autorización de medidor y membresía. Se rechazan campos desconocidos, unidades incorrectas, valores no finitos, rangos no plausibles y horas fuera de la ventana admitida. El UUID público usado en el transporte difiere del identificador interno de la fila de dispositivo.

\subsection{Cálculo por contadores y cobertura temporal}
Si $C(t)$ representa energía acumulada, el incremento observado es:
\begin{equation}
\Delta E=C(t_1)-C(t_0), \qquad t_1>t_0.
\end{equation}
La primera observación establece una línea base. Un mensaje repetido no genera un incremento adicional. Si el contador disminuye, se trata como reinicio salvo que exista una especificación de rollover autorizada y se cumplan sus verificaciones. Un hueco no se convierte automáticamente en energía estimada.

Los intervalos habituales son de 900 segundos, alineados en UTC. La cadencia normal es de 15 segundos y el cierre espera dos minutos de gracia. La elegibilidad requiere observaciones procesadas en los dos límites exactos, además de autorización vigente y ausencia de evidencia inválida. El planificador incluye límites contables aun cuando la cadencia configurada no los divide; los cortes perdidos no se rellenan retroactivamente.

Una muestra inicial parcial o un periodo posterior a una desconexión puede resultar incompleto o rechazado. Esta es una exclusión contable, no una instrucción para cambiar manualmente el estado a finalizado.

\subsection{Calidad y trazabilidad}
\begin{longtable}{L{.26\linewidth}L{.64\linewidth}}
\caption{Interpretación de estados energéticos y de telemetría.}\\
\toprule Estado & Significado operativo \\ \midrule\endfirsthead
\toprule Estado & Significado operativo \\ \midrule\endhead
\code{valid} & Sobre aceptado y pendiente de procesamiento energético.\\
\code{PROCESSED} & Mensaje procesado; no prueba por sí solo que todo el intervalo sea elegible.\\
\code{COUNTER_GAP}, \code{COUNTER_RESET} & Se requiere una nueva base; no se acredita como energía el salto problemático.\\
\code{OPEN} & Intervalo aún en cálculo o pendiente de cierre y evidencia.\\
\code{INCOMPLETE} & No hay cobertura de límites suficiente para la política actual.\\
\code{REJECTED}, \code{DISPUTED} & Evidencia o autorización no aceptable, o contradicción que exige revisión.\\
\code{FINALIZED} en intervalo & El cálculo energético terminó; no equivale a lote aprobado ni token emitido.\\
\bottomrule
\end{longtable}
Las evaluaciones de calidad conservan motivo, digest y ocurrencias. Un mensaje inválido no debe refrescar la señal de última telemetría autenticada. En redes no confiables se dispone de provisión de ACL por tópico y overlay TLS; la configuración local anónima no se considera segura para exposición pública.

\section{Modelo de datos e integridad}
\subsection{Fuente del esquema y estado revisado}
SQLModel y Alembic son la fuente operativa del esquema. Durante esta revisión se registraron \textbf{44 tablas} y se comprobó la cabecera \code{20261112_local_recovery}. Los archivos SQL históricos y el snapshot OpenAPI v2 no se utilizan para instalar una estructura paralela. Las cifras se refieren a tablas registradas por el código inspeccionado, no a una certificación de que cualquier base del usuario esté migrada correctamente \citep{database,models,migration}.

La imagen Docker incluye TimescaleDB, pero las migraciones vigentes no crean hypertables ni retención automática. Por tanto, se documenta una persistencia PostgreSQL relacional con posibilidad futura de diseño temporal, sin atribuirle particionamiento no implementado.

\subsection{Dominios principales}
\begin{longtable}{L{.25\linewidth}L{.65\linewidth}}
\caption{Agrupación de las tablas vigentes.}\label{tab:dominios}\\
\toprule Dominio & Tablas y relaciones destacadas \\ \midrule\endfirsthead
\toprule Dominio & Tablas y relaciones destacadas \\ \midrule\endhead
Identidad & \code{users}, \code{wallets}, \code{wallet_nonces}. La clave de usuario es \code{users.uid}; las wallets se relacionan con ella.\\
Infraestructura & \code{homes}, \code{devices}, \code{device_credentials}. El hogar tiene propietario y el dispositivo dos identidades: interna y pública.\\
Comunidad & \code{energy_communities}, \code{community_memberships}, \code{community_nodes}, \code{meter_authorizations}. Delimitan participación y vínculo del medidor con el nodo.\\
Telemetría & \code{telemetry_messages}, \code{telemetry_assessments}; lecturas tipadas en \code{smart_meter}, \code{inverter}, \code{battery}, \code{weather}.\\
Contabilidad & \code{device_counter_snapshots}, \code{energy_power_snapshots}, \code{energy_intervals}, \code{energy_credits}, \code{credit_movements}.\\
Protocolo nativo & \code{native_settlement_policies}, \code{native_settlement_batches}, \code{native_settlement_claims}, \code{native_interval_allocations}, \code{community_native_deployments}.\\
Ejecución y proyección & \code{outbox_events}, \code{community_chain_operations}, \code{native_indexed_events}, \code{native_sync_cursors}, \code{native_wallet_transactions}.\\
Auditoría y recuperación & \code{local_chain_recoveries}, \code{community_energy_ledger}, \code{audit_events}.\\
Historia retenida & Tablas de órdenes, trades, redenciones, despliegues ERC-20 y obligaciones previas; su existencia no las convierte en flujo nativo activo.\\
\bottomrule
\end{longtable}

\begin{figure}[htbp]
\caption{Relaciones principales del modelo de datos}\label{fig:datos}

\centering
\begin{tikzpicture}[node distance=.85cm and 1.05cm,box/.append style={text width=3.7cm}]
\node[box] (u) {Usuario\\\code{users.uid}};
\node[box,right=of u] (h) {Hogar\\\code{homes.owner_id}};
\node[box,right=of h] (d) {Dispositivo\\\code{devices.home_id}};
\node[box,below=of u] (c) {Comunidad y membresía};
\node[box,below=of h] (n) {Nodo y autorización\\del medidor};
\node[box,below=of d] (t) {Mensajes y\\lecturas tipadas};
\node[box,below=of t] (i) {Intervalos y créditos};
\node[box,left=of i] (b) {Política y lote};
\node[box,left=of b] (a) {Claims y asignaciones};
\draw[arr] (u)--node[labelnode,above]{1:N}(h);
\draw[arr] (h)--node[labelnode,above]{1:N}(d);
\draw[arr] (u)--(c);
\draw[arr] (c)--(n);
\draw[arr] (h)--(n);
\draw[arr] (n)--(d);
\draw[arr] (d)--(t);
\draw[arr] (t)--(i);
\draw[arr] (c.south) to[out=-70,in=170] (b.west);
\draw[arr] (b)--(a);
\draw[arr] (i.south) to[out=-120,in=-70] (a.south);
\end{tikzpicture}
\par\noindent\small\textit{Nota.} Elaboración propia a partir de los modelos SQLModel. Las flechas de evidencia indican trazabilidad lógica; no todas corresponden a una clave foránea directa \citep{models}.

\end{figure}

\subsection{Tipos y restricciones}
La energía se representa en Wh enteros. Las obligaciones y precios nativos utilizan \code{NUMERIC(78,0)} para conservar precisión de enteros compatibles con uint256. En la API se transmiten como cadenas decimales. SUN utiliza 18 decimales y los tiempos de cálculo y eventos son conscientes de zona horaria.

Se aplican restricciones de unicidad para correo y dirección de wallet, secuencia y mensaje por dispositivo, intervalo por hogar, crédito por intervalo, lote por comunidad y periodo, claim por lote/tipo/nodo y operación firmada por outbox. Los eventos indexados incluyen identidad de bloque para conservar versiones huérfanas sin tratarlas como canónicas.

Los bloqueos y las transacciones complementan estas restricciones. Una comprobación previa en Python no evita por sí sola que dos trabajadores concurrentes firmen la misma obligación. El ledger comunitario añade un trigger que rechaza actualización y borrado; un administrador de base de datos sigue teniendo capacidades superiores de mantenimiento, por lo que no se afirma inviolabilidad.

\subsection{Separación entre evidencia y proyección}
La base conserva la telemetría, su calidad y las referencias de cálculo. En cambio, saldos, entregas y retiros confirmados tienen autoridad en los contratos. Las tablas nativas contienen proyecciones que deben actualizarse desde recibos y eventos canónicos. Insertar una fila \code{ISSUED} o \code{CONFIRMED} sin transacción no constituiría emisión ni pago.

Las hojas y pruebas Merkle se conservan junto con las cantidades, nodo, medidor, beneficiario y política. No deben regenerarse silenciosamente a partir de estado mutable después de aprobar un lote.

\subsection{Semillas y recuperación}
El seed SQLModel crea cinco hogares y quince dispositivos de ejemplo y respeta identidades, contraseñas y revocaciones existentes. El bootstrap gobernado es una fase diferente: despliega contratos y autoriza participantes mediante transacciones reales locales.

La migración de recuperación añade \code{archived}, origen de génesis y un registro \code{local_chain_recoveries}; no elimina la historia. La nueva marca \code{accounting_start} impide volver a emitir desde intervalos anteriores al reemplazo de cadena. Un registro archivado en \code{RECONCILIATION_REQUIRED} conserva evidencia de una cadena perdida: no es una promesa de que esperar restaurará sus fondos.

\section{Contratos, gobernanza y reglas económicas}
\subsection{Contratos del protocolo activo}
\begin{longtable}{L{.28\linewidth}L{.62\linewidth}}
\caption{Responsabilidades de los cinco contratos activos.}\label{tab:contratos}\\
\toprule Contrato & Responsabilidad \\ \midrule\endfirsthead
\toprule Contrato & Responsabilidad \\ \midrule\endhead
\code{TimelockController} & Programación y ejecución con retraso y roles de gobernanza.\\
\code{CommunityRegistry} & Ciclo de vida, membresías, nodos, medidores, precio y vínculos con los contratos comunitarios.\\
\code{SUNToken} & ERC-20 con emisión restringida, 18 decimales, pausado, quema y anclajes contables.\\
\code{CommunitySettlement} & Lotes, revisión, finalización, prueba de inyección, emisión acotada y retiro por consumo.\\
\code{NativeSettlementVault} & Depósitos de ETH y SUN, reservas, entrega atómica y retiros por el propietario.\\
\bottomrule
\end{longtable}
Los contratos utilizan Solidity 0.8.24 y componentes OpenZeppelin. Los ABI nativos compartidos entre backend y frontend se sincronizan mediante un script específico. La administración se asigna al contrato timelock. El minter de SUN se concede a CommunitySettlement y se revoca al otorgante inicial una vez autorizado el protocolo \citep{blockchain,contracts}.

\subsection{Ciclo de vida y gobernanza}
La comunidad pasa por propuesta, revisión, aprobación y activación, con transiciones autorizadas para rechazo, suspensión o archivo. La membresía también tiene estados y no se obtiene por conectar una wallet. El operador comunitario y los roles de registrador de medidores, oráculo, liquidación y pausa son permisos diferentes.

Una acción del timelock se \textbf{programa} y después se \textbf{ejecuta} al cumplirse el retraso. El entorno local utiliza demora cero y conserva las dos transacciones de programación y ejecución. El lote también tiene un periodo de revisión antes de su finalización.

El precio por kWh se establece mientras la comunidad está aprobada y queda fijado al activarla. Ese precio se incorpora a las asignaciones del lote.

\subsection{Asignación de cantidades}
Sea $S$ la inyección neta elegible total y $D$ la demanda neta elegible. La cantidad asignable cumple:
\begin{equation}
A=\min(S,D),\quad 0\leq A\leq S,\quad A\leq D.
\end{equation}
La política compensa importación y exportación del intervalo y la autooferta de varios nodos de una misma wallet. La política permite aplicar un factor de pérdidas a la oferta elegible.

Para repartir $A$ entre demandas $d_j$, la base entera es:
\begin{equation}
a_j^{(0)}=\left\lfloor\frac{A d_j}{\sum_k d_k}\right\rfloor.
\end{equation}
Las unidades restantes se distribuyen por mayores residuos y desempate estable de identidad. Se realiza un reparto análogo de oferta y un emparejamiento determinista productor-consumidor. El reparto se realiza entre participantes distintos. El uso de Wh enteros evita residuales energéticos de punto flotante \citep{allocation,settlement}.

\subsection{Lotes y compromisos de evidencia}
El oráculo presenta intervalo, política, raíz Merkle, totales y asignación en un lote. El estado inicial en cadena es \code{CHALLENGEABLE}. Después de la revisión, la gobernanza puede finalizarlo o rechazarlo. Solo un lote canónico puede quedar finalizado por comunidad e intervalo bajo las reglas actuales.

Una prueba Merkle demuestra pertenencia a la raíz comprometida. Las hojas incluyen clase de operación, lote, claim o asignación, nodos, medidores, beneficiarios y cantidades. La doble aplicación de keccak y el orden de pares coinciden entre ABI, backend y contratos. La raíz permite comprobar que una prueba corresponde al conjunto aprobado.

El contrato permite sustituir un lote finalizado que conserve las condiciones de uso requeridas.

\subsection{Emisión, unidades y precio}
La equivalencia contable es:
\begin{equation}
1\;\text{SUN}=1\;\text{kWh},\qquad
1\;\text{Wh}=10^{15}\;\text{unidades base SUN}.
\end{equation}
La finalización del lote habilita la emisión. El trabajador tramita la transacción y consulta su confirmación. La emisión agregada no puede exceder la inyección aprobada y cada claim es único. La emisión precede a la financiación y liquidación de las asignaciones.

El precio $p$ procede de la política comunitaria. Para una asignación de $w$ Wh, el coste en wei es:
\begin{equation}
c=\left\lceil\frac{w\,p}{1000}\right\rceil.
\label{eq:coste}
\end{equation}
En el ejemplo, $p=10^{15}$ wei/kWh, equivalente a 0,001 ETH/kWh. Así, 108 Wh corresponden a 0,108 SUN y 0,000108 ETH. Fraccionar cantidades en varias parejas puede añadir hasta un wei de redondeo por pareja frente a una única cuenta conjunta; el mismo importe se acredita y se adeuda, sin residuo monetario sin propietario.

\subsection{Depósitos, reserva y entrega}
El productor primero autoriza gasto ERC-20 al vault y después deposita SUN. El consumidor deposita ETH nativo y establece un límite de precio. El participante firma el aporte de ETH desde su wallet. Una reserva exige prueba de asignación, cantidad, precio, ETH y SUN suficientes; si falta algún respaldo, la transacción revierte completa.

La ejecución de una reserva válida entrega SUN al consumidor y acredita ingresos retirables del productor en la misma transacción. El ingreso se acredita durante la liquidación y se cobra con \code{withdrawEarnings}. La cancelación de una reserva libera los activos correspondientes.

\begin{figure}[htbp]
\caption{Flujo de activos durante la liquidación}\label{fig:activos}

\centering
\begin{tikzpicture}[node distance=.85cm and 1.2cm,box/.append style={text width=3.6cm}]
\node[box] (prod) {Productor\\SUN en wallet};
\node[box,right=of prod] (v) {Vault\\SUN y ETH reservados};
\node[box,right=of v] (cons) {Consumidor\\ETH en wallet};
\node[box,below=1.2cm of prod] (earn) {Ingresos acreditados\\y retiro de ETH};
\node[box,below=1.2cm of v] (exec) {Ejecución atómica\\de la asignación};
\node[box,below=1.2cm of cons] (burn) {SUN recibidos\\y retiro por consumo};
\draw[arr] (prod)--node[labelnode,above]{autorizar y depositar}(v);
\draw[arr] (cons)--node[labelnode,above]{depositar ETH}(v);
\draw[arr] (v)--(exec);
\draw[arr] (exec)--(earn);
\draw[arr] (exec)--(burn);
\end{tikzpicture}
\par\noindent\small\textit{Nota.} Elaboración propia a partir del protocolo implementado. La quema posterior de SUN registra el consumo y no añade un pago ETH \citep{settlement}.

\end{figure}

\subsection{Retiro por consumo e invariantes}
El consumidor aprueba un allowance diferente, dirigido a CommunitySettlement, y presenta una prueba de consumo. El retiro exige lote finalizado, participante/nodo/medidor autorizados y entrega financiera completa de la asignación comprometida para ese nodo. Se queman los SUN y se marca el consumo como utilizado. La prueba vincula los SUN utilizados con el consumo entregado.

Los invariantes principales de respaldo son:
\begin{align}
\operatorname{ETH}(vault)&\geq \sum_u disponible_u+\sum_u reservado_u+\sum_u ingresos_u,\\
\operatorname{SUN}(vault)&\geq \sum_u SUNdisponible_u+\sum_u SUNreservado_u.
\end{align}
Las autorizaciones, la unicidad de claims, el límite de emisión y las protecciones contra reentrancia complementan estas igualdades de solvencia. El contrato de SUN permite transferencias fungibles; la trazabilidad de inyección y consumo se conserva en los registros de emisión, asignación y retiro, no en la identidad individual de cada unidad del token.

\section{Funcionalidades de la plataforma}
\subsection{Acceso y wallet}
La aplicación permite crear una cuenta e iniciar sesión. La wallet se conecta por separado y se vincula al usuario mediante la firma de un mensaje con un código de un solo uso. El servidor comprueba esa firma y registra la dirección del participante.

Las acciones sobre fondos utilizan la wallet vinculada al usuario. La interfaz muestra la dirección, la red y el estado de verificación. Los permisos de gobernanza se consultan en los contratos.

\subsection{Pantallas y operaciones}
\begin{longtable}{L{.23\linewidth}L{.67\linewidth}}
\caption{Funciones de las pantallas de la aplicación}\label{tab:funcionalidades}\\
\toprule Pantalla & Funciones \\ \midrule\endfirsthead
\toprule Pantalla & Funciones \\ \midrule\endhead
Inicio & Consultar generación, consumo, exportación, SUN, obligaciones e ingresos.\\
Hogares & Registrar, consultar y editar ubicación, zona horaria, capacidad solar y batería.\\
Dispositivos & Registrar equipos, seleccionarlos por hogar y consultar lecturas, salud y calidad.\\
Energía & Filtrar por hogar y periodo, revisar gráficas y consultar intervalos.\\
Comunidad & Consultar reglas y nodos, presentar propuestas y solicitar membresía.\\
Créditos SUN & Consultar cantidades, estados de emisión y relación con el intervalo de origen.\\
Liquidación & Revisar lotes, asignaciones, financiación y pruebas; consultar las tareas del operador.\\
Gobernanza & Preparar, programar y ejecutar acciones de la cuenta autorizada.\\
Depósitos y retiros & Depositar ETH, autorizar y depositar SUN, retirar activos disponibles y cobrar ingresos.\\
Compensación & Seleccionar consumo entregado, autorizar SUN y registrar su quema.\\
Transacciones & Consultar la actividad de la sesión y el registro de operaciones del servidor.\\
Configuración & Consultar el perfil, el estado de servicios y la zona horaria utilizada.\\
\bottomrule
\end{longtable}

\subsection{Operaciones del productor}
El productor recibe SUN por la inyección aprobada. Para usarlos en la liquidación, autoriza una cantidad al contrato de custodia y luego ejecuta el depósito. La aplicación muestra los SUN en wallet, los disponibles en depósito y los reservados.

Cuando se ejecuta una asignación, el contrato acredita ETH al productor. Este ingreso aparece como reclamable y se envía a su wallet al firmar el retiro. También puede recuperar los SUN depositados que estén disponibles.

\subsection{Operaciones del consumidor}
El consumidor consulta la energía que tiene asignada y su coste. Deposita ETH con un precio máximo permitido; el contrato reserva los fondos correspondientes a cada asignación. Al liquidarse, recibe SUN.

En Compensación selecciona el consumo entregado, autoriza el uso de los SUN y firma la operación que los quema. El registro resultante vincula esa cantidad con el consumo aprobado. El consumidor también puede retirar ETH disponible de su depósito.

\subsection{Gobernanza}
La cuenta de gobernanza administra las autorizaciones de la comunidad y la finalización de lotes. La aplicación prepara la acción y la wallet firma su programación y ejecución. La finalización confirmada habilita el procesamiento de las emisiones correspondientes.

\subsection{Revisión y firma}
Antes de una operación, la aplicación presenta red, cuenta, contrato, valor y argumentos. La firma se realiza en MetaMask. Autorizar un gasto y transferir los activos son acciones separadas: \code{approve_sun} autoriza el depósito y \code{escrow_sun} lo ejecuta. Para el consumo, \code{approve_retirement} autoriza la quema y \code{retire} la registra.

\subsection{Estados de la operación}
El usuario puede seguir cuatro estados de negocio: intervalo calculado, lote aprobado, crédito emitido y asignación liquidada. El estado \code{FINALIZED} pertenece al cálculo o al lote según la pantalla; \code{ISSUED} identifica un crédito emitido y \code{SETTLED} una asignación liquidada.

La financiación se consulta por separado. El estado \textit{Con fondos} indica que hay recursos suficientes o una reserva constituida. La entrega se confirma mediante el estado de liquidación.

\subsection{Unidades y saldos}
Un SUN representa un kWh de energía elegible aprobada. El token utiliza dieciocho decimales para registrar fracciones. Por ejemplo, una asignación de 0,108 kWh corresponde a 0,108 SUN.

Los saldos de wallet, depósito, reserva e ingreso corresponden a ubicaciones y usos distintos de los activos. La consulta separada permite conocer qué cantidad puede depositarse, retirarse o utilizarse en una asignación.

\chapter{Despliegue, continuidad y seguridad}
\section{Ciclo de vida Docker}
El punto de entrada operativo es \code{scripts/stack.py}, invocado por Make. Todos los comandos comparten el proyecto original y los archivos Compose canónicos. El proyecto selecciona los volúmenes persistentes; cambiarlo para evitar un error puede conectar otra instalación y ocultar los registros antiguos.

\begin{longtable}{L{.23\linewidth}L{.67\linewidth}}
\caption{Comandos mantenidos y efecto sobre la instalación.}\label{tab:comandos}\\
\toprule Comando & Efecto \\ \midrule\endfirsthead
\toprule Comando & Efecto \\ \midrule\endhead
\code{make setup} / \code{make up} & Construyen lo necesario, aplican migraciones, comisionan o verifican el protocolo local y levantan la aplicación conservando identidades existentes.\\
\code{make simulate} & Levanta la aplicación completa y el simulador registrado; no requiere dejar una app parcial iniciada por otro archivo Compose.\\
\code{make doctor} / \code{make ps} & Inspeccionan dependencias, disco, despliegue, readiness y contenedores.\\
\code{make backup} & Exporta un archivo SQL privado y verifica su estructura con pg\_restore.\\
\code{make down} & Detiene servicios y conserva volúmenes.\\
\code{make recover} & Recupera una instalación local sobre el mismo PostgreSQL; si es necesario reemplaza cadena y archiva historia económica.\\
\code{make clean} & Elimina únicamente artefactos de compilación de contratos, no registros ni volúmenes.\\
\code{make demo-scenario} & Ejecuta o reanuda la prueba finita MQTT/API/EVM de 5/3 kWh.\\
\code{make blockchain-check} & Inserta telemetría validada directamente y verifica el mismo escenario; no prueba el transporte MQTT.\\
\bottomrule
\end{longtable}
Los servicios incluyen PostgreSQL/TimescaleDB, Redis, Mosquitto, Anvil, inicialización de estado, migración, API, worker, suscriptor, frontend y simulador opcional. API, worker y suscriptor esperan migración correcta y dependencias saludables. Los puertos públicos locales se limitan a loopback: 8080 para frontend, 8000 para API, 8545 para RPC y 1883 para MQTT.\cite{deployment,start}

\section{Persistencia EVM y recuperación}
La pérdida de Anvil fue una limitación operativa relevante: podían conservarse usuarios y lecturas mientras las direcciones almacenadas ya no contenían contratos. El error \code{BadFunctionCallOutput} y la ausencia de bytecode no se resuelven acreditando saldos en SQL.

El supervisor actual mantiene snapshots validados de hasta 128 MiB, con escritura en archivo temporal, fsync y reemplazo atómico; conserva un checkpoint anterior completo. Obtiene snapshots cada 30 segundos y durante la parada ordenada. Limita estados históricos y evita el crecimiento no controlado de logs mediante rotación. Si el snapshot actual falla, intenta el previo. Si ambos son inválidos, el arranque se detiene; no continúa silenciosamente en una cadena vacía.

Un corte abrupto puede perder operaciones posteriores al último checkpoint. Por eso conservar un archivo válido no elimina la necesidad de conciliación. La identidad de génesis se preserva y se comprueba; los relojes de minería se ajustan sobre el último estado cargado para no quedar atascados cerca de la fecha inicial.

Cuando la cadena es irrecuperable, \code{make recover}:
\begin{enumerate}
\item Detiene procesos que escriben y verifica acceso al mismo proyecto y base.
\item Crea un respaldo SQL privado verificable y conserva configuración y journals.
\item Archiva snapshots de Anvil y el historial financiero relacionado con la cadena perdida.
\item Registra una recuperación con referencias antiguas/nuevas y conserva las filas.
\item Despliega contratos locales nuevos y reafirma identidades y autorizaciones originales.
\item Fija un corte contable para evitar nueva emisión de energía histórica.
\item Conserva un marcador de progreso, checkpoint y configuración para poder reanudar una interrupción.
\end{enumerate}
Restaurar SQL no reconstruye almacenamiento EVM, recibos ni saldos. Los lotes archivados mantienen evidencia histórica; el usuario no debe esperar que una conciliación automática resucite una cadena borrada.

\section{Corrección de credenciales PostgreSQL}
La configuración previa forzaba \code{solar:solar} dentro del backend de Compose. Un volumen creado con otra contraseña podía aparecer saludable por pg\_isready, pero la migración fallaba con \code{InvalidPasswordError}. Se corrigió el despliegue para conservar la URL original y adaptar únicamente host y puerto a \code{timescaledb:5432} dentro del contenedor.

El módulo \code{src.core.docker_database} utiliza el parser de URL de SQLAlchemy, conserva credenciales codificadas, base y parámetros, y evita imprimir la URL en errores. El entrypoint exporta el resultado antes de la migración o la app. \code{POSTGRES_PASSWORD} inicializa volúmenes nuevos; editarlo no cambia la contraseña de un rol existente. La corrección de configuración no rota credenciales ni modifica registros de negocio.\cite{dbfix}

\section{Salud y diagnóstico}
\code{/health} comprueba vivacidad de API/base; \code{/ready} agrega Redis, red y código de registro, además de actividad de worker y suscriptor. Un estado ready indica disponibilidad actual, no que las lecturas sean válidas ni que un lote concreto se haya liquidado.

Las métricas del suscriptor incluyen persistidos, duplicados, rechazados y cola. Un máximo de hora de medición sin cambios entre dos consultas requiere compararlo con el reloj y los eventos recientes, no declarar inmediatamente detenido el simulador. Un mensaje rechazado puede ser antiguo y un contador de ventana móvil puede disminuir sin caída del proceso. El diagnóstico combina logs acotados, calidad por dispositivo, límites de intervalos y estados de lote/financiación.

\section{Seguridad implementada y límites}
Las contraseñas se almacenan con bcrypt; los JWT usan revocación de JTI en Redis y las verificaciones de wallet tienen nonce único. Las credenciales de medidor se comparan con digests y se excluyen del almacenamiento de evidencia. Los contenedores de backend operan sin root. Las firmas, nonces y génesis se registran antes de transmitir operaciones de servicio.

Persisten dependencias de confianza y necesidades de endurecimiento: tokens legibles por el navegador, refresh JWT sin rotación opaca, logout limitado al token presentado, correo de recuperación no desplegado como servicio canónico, gestión de secretos y alta disponibilidad pendientes. Los brokers anónimos locales y las claves de fixtures no son apropiados para un servicio público. Las operaciones sensibles de contratos usan permisos y protecciones de reentrancia, pero no existe una auditoría externa que certifique ausencia de vulnerabilidades.\cite{security}

\section{CI y reproducibilidad}
El workflow \code{.github/workflows/verify.yml} configura instalación de dependencias, suites de contratos/backend/IoT/frontend, chequeo de ABI, stack, escenario, reinicio del worker y pruebas de navegador. Su existencia demuestra la automatización configurada, no el éxito de una ejecución remota concreta; para afirmarlo se necesitaría el registro del job correspondiente.

Las dependencias se fijan con instaladores y lockfiles. La compilación y el despliegue emplean versiones distintas de herramientas de host y del nodo Docker según sus guías; no se supone que todos los ejecutables Foundry tengan la misma versión.

\section{Verificación y resultados alcanzados}
\subsection{Escenario financiero versionado}
La evidencia principal es \code{docs/evidence/local-native-scenario.json}, con estado \code{VERIFIED}. El archivo identifica comunidad, lote, intervalo, raíz, estados de cuenta y hashes de operaciones. Su intervalo retenido es del 11 de noviembre de 2026, de 01:30 a 01:45 UTC; las horas locales deben convertirse antes de compararlas con una captura. No se afirma que esos hashes sigan disponibles en un Anvil reemplazado posteriormente ni que se puedan buscar en un explorador público \citep{scenario,verification}.

\begin{longtable}{L{.53\linewidth}L{.37\linewidth}}
\caption{Resultados del escenario local de 5/3 kWh.}\label{tab:escenario}\\
\toprule Magnitud & Resultado confirmado en el artefacto \\ \midrule\endfirsthead
\toprule Magnitud & Resultado confirmado en el artefacto \\ \midrule\endhead
Inyección elegible / demanda & 5 kWh / 3 kWh\\
Asignación / inyección no asignada & 3 kWh / 2 kWh\\
Precio aplicado & 0,001 ETH de prueba por kWh\\
Depósito inicial del consumidor & 0,01 ETH\\
Pago de las asignaciones & 0,003 ETH\\
SUN emitidos / entregados / retirados & 5 / 3 / 3 SUN\\
ETH disponible final del consumidor & 0,007 ETH\\
Ingreso retirado del productor & 0,003 ETH\\
Oferta SUN restante & 2 SUN\\
Pasivos finales del vault & 0,007 ETH\\
Reservas e ingresos pendientes al finalizar & 0 ETH\\
\bottomrule
\end{longtable}
Las cuentas cierran: 0,01 ETH aportados menos 0,003 ETH retirados por el productor dejan 0,007 ETH respaldando disponibilidad del consumidor; 5 SUN emitidos menos 3 SUN quemados dejan 2 SUN. La energía no asignada no genera por sí misma ingreso: quedan créditos del productor, no otro pago automático.

El escenario verifica autenticación, deduplicación y reanudación mediante registro de operaciones. Los scripts conservan hashes antes de esperar recibos para no repetir depósitos, retiros o quema. \code{blockchain-check} comparte el mismo escenario con \code{demo-scenario}: reejecutarlo después de completar no crea otro pago. Tras una recuperación se utilizan bases de contadores y un intervalo posterior al corte; no se borran journals para forzar reutilización de créditos.

\subsection{Resultados de pruebas por campaña}
\begin{longtable}{L{.24\linewidth}L{.26\linewidth}L{.39\linewidth}}
\caption{Evidencias automatizadas separadas por alcance y origen.}\label{tab:pruebas}\\
\toprule Campaña / fuente & Resultado retenido & Interpretación y límites \\ \midrule\endfirsthead
\toprule Campaña / fuente & Resultado retenido & Interpretación y límites \\ \midrule\endhead
Integración documentada en \code{e061ab0} & 38 contratos; 48 backend; 86 IoT; 35 frontend; 11 navegador. & Cifras históricas de \code{docs/verification.md}; no son el total actual ni se vuelven a ejecutar para esta redacción.\\
Revisión del despliegue \code{2ef7536} & Backend: 80 pruebas; IoT: 95 pruebas, 88,18\% de cobertura. & Registros de campañas locales posteriores; no sumables a los anteriores como pruebas distintas.\\
Conservación SQL de recuperación & 44 tablas con claves conservadas; contraseñas, propiedad y payloads históricos preservados. & Comparación local con respaldo de inspección; no certifica restauración de cualquier desastre.\\
Escenario posterior a recuperación & Ciclo de 5/3 kWh verificado nuevamente. & Ejecución local retenida en logs de trabajo; conserva el mismo alcance sintético y no público.\\
Inserción validada \code{46d5da8} & 26 pruebas seleccionadas de pipeline y ciclo de vida. & Incluye PostgreSQL/Anvil aislados y rechazo de credenciales/datos alterados; no equivale a toda la suite.\\
Comprobación de este informe & 32 pruebas de ciclo de vida aprobadas. & Suite seleccionada ejecutada durante la revisión; no se lanzó una campaña financiera nueva.\\
\bottomrule
\end{longtable}
La cobertura de IoT es una medición de líneas instrumentadas de esa campaña; no es cobertura del sistema completo, garantía de corrección o índice de seguridad. Los manifiestos de evidencia adjuntos resumen los resultados sin incluir logs privados ni credenciales.

\subsection{Comprobaciones estructurales actuales}
En la revisión base se importó el registro de modelos, se generó OpenAPI desde la aplicación y se consultó Alembic sin conectar una nueva base de negocio. Los resultados son 44 tablas, 71 paths y 80 operaciones HTTP, con cabecera \code{20261112_local_recovery}. Estos números incluyen estructuras y endpoints históricos retenidos: no son 80 funcionalidades nativas distintas expuestas al usuario.



\subsection{Matriz de cumplimiento}
\begin{longtable}{L{.29\linewidth}L{.18\linewidth}L{.42\linewidth}}
\caption{Estado de objetivos frente a la evidencia revisada.}\label{tab:cumplimiento}\\
\toprule Capacidad & Estado & Sustento o pendiente \\ \midrule\endfirsthead
\toprule Capacidad & Estado & Sustento o pendiente \\ \midrule\endhead
Telemetría autenticada y auditable & Logrado localmente & Validación v2, deduplicación, fuente y calidad; pruebas con broker y pipeline.\\
Contabilidad por intervalos & Logrado localmente & Contadores, límites UTC, gracia y exclusión de evidencia defectuosa.\\
Gobernanza de lotes & Implementado y verificado & Timelock, prepare de acciones, firmas y finalización antes de emisión.\\
Emisión acotada y retiro por consumo & Verificado en escenario & Claims únicos, prueba, límite agregado y quema de consumo entregado.\\
Liquidación financiada & Verificado en escenario & Custodia doble, reserva, entrega SUN y retirada de ingresos ETH.\\
Recuperación sin eliminar SQL & Verificado en ensayo local & Comparación de tablas/identidades, archivo y corte contable.\\
Interfaz operativa en español & Implementada y probada localmente & Tests de componentes/navegador y observaciones manuales; no auditoría completa de accesibilidad.\\
Medición física calibrada & Pendiente & No hay evidencia de instalación y contraste metrológico.\\
Sepolia y MetaMask público & Pendiente de verificación & Tooling disponible; sin aceptación pública registrada.\\
Facturación y operación comercial & Fuera del resultado actual & Sin integración de factura ni validación regulatoria de venta.\\
\bottomrule
\end{longtable}

\subsection{Lectura de los resultados}
El aporte principal es cerrar un ciclo de software con respaldo verificable en la misma red de prueba: datos, aprobación, emisión, fondos, entrega y salidas. La evidencia respalda integridad funcional en escenarios ensayados y cumplimiento de las reglas de autorización y respaldo. No respalda una estimación de ahorro, reducción de emisiones, capacidad máxima de usuarios, latencia p95, transacciones por segundo, disponibilidad anual o coste de gas en producción. Esas métricas requieren experimentos propios, condiciones de carga y mediciones no presentes.

\chapter{Conclusiones, limitaciones y continuidad}
\section{Conclusiones del proceso}
\begin{enumerate}
\item Se obtuvo una integración local que conserva el vínculo entre medición, autorización, crédito y pago. El resultado no depende de marcar manualmente como pagada una obligación SQL.
\item Separar identidad de aplicación, propiedad de wallet y roles en cadena fue necesario para explicar y proteger la operación. La experiencia de demostración evidenció que esta separación debe comunicarse mejor al usuario.
\item El cálculo conservador por contadores y límites observados impide acreditar huecos como energía. La contrapartida es que una interrupción puede dejar intervalos sin elegibilidad y requerir esperar un periodo completamente observado.
\item La gobernanza finaliza la evidencia; la emisión posterior, la financiación y la entrega son etapas diferentes. Esta separación explica estados que de otro modo podrían interpretarse como contradicciones del sistema.
\item La persistencia y la génesis forman parte de la corrección económica. Conservar PostgreSQL sin el estado de Anvil no recupera saldos EVM. La recuperación debe preservar historia y evitar reemisión, no ocultar el fallo con una base nueva.
\item La documentación consolidada y el escenario verificable permiten reproducir el alcance local. Las cifras de suites y cobertura deben mantenerse ligadas a sus campañas y revisiones.
\end{enumerate}

\section{Limitaciones materiales}
La medición demostrada es simulada. Los contratos solo conocen pruebas y compromisos aportados por actores autorizados; un medidor legítimo que mienta no queda descubierto automáticamente por Merkle. No existe auditoría independiente de contratos, infraestructura o autenticación. Los límites de escalabilidad y los costes públicos no se han medido.

La política actual tiene precio fijo aprobado antes de activación, sin mercado abierto de órdenes nativo ni formación dinámica de precio. La corrección de derechos ya pagados no es automática. Existen carencias de perfil editable, descubrimiento público, manifiestos completos de evidencia, historias de eventos SUN y recuperación de contraseña desplegada.

No se han demostrado instalaciones físicas, acuerdos con un operador de red, validez regulatoria para venta o modificación de facturas. No se propone que SUN sustituya moneda de curso legal ni que ETH de prueba tenga valor comercial.

\section{Trabajo futuro priorizado}
\begin{longtable}{L{.14\linewidth}L{.44\linewidth}L{.31\linewidth}}
\caption{Continuidad propuesta con criterios verificables.}\\
\toprule Prioridad & Trabajo propuesto & Criterio de cierre \\ \midrule\endfirsthead
\toprule Prioridad & Trabajo propuesto & Criterio de cierre \\ \midrule\endhead
Alta & Mejorar selección y advertencias de wallet; normalizar unidades y nombres de acciones. & Recorrido multiusuario sin confusión entre cuentas; validación con usuarios y pruebas de interfaz.\\
Alta & Auditoría externa, modelo de amenazas y gestión de claves de oráculo/gobernanza. & Hallazgos revisados y controles verificados antes de exposición pública.\\
Alta & Instalar medidores físicos con calibración y commissioning verificable. & Comparación contra referencias metrológicas y trazabilidad del dispositivo.\\
Alta & Ensayar restauración conjunta SQL/EVM y fallos de infraestructura. & Ensayos medidos, pérdidas máximas registradas y runbook de recuperación.\\
Media & Completar despliegue Sepolia con fuentes admitidas y autoridades separadas. & Contratos, permisos, recibos y prueba de aceptación pública conservados.\\
Media & Medir carga, latencia, backlog, almacenamiento y costes. & Dataset y protocolo de benchmark reproducibles; no extrapolar desde el demo.\\
Media & Revisar claves temporales para hypertables y políticas de retención. & Migración compatible con referencias y unicidad, verificada sobre datos poblados.\\
Posterior & Integración de facturación y análisis regulatorio colombiano. & Requisitos y validaciones externas documentados antes de afirmar compensación comercial.\\
\bottomrule
\end{longtable}

\section{Resultado final alcanzado}
El resultado es un \textbf{prototipo funcional y demostrable de contabilidad y liquidación energética comunitaria en red local}. La medición respalda una cantidad elegible; la gobernanza autoriza el compromiso; la emisión representa el derecho; la custodia asegura recursos; y la ejecución confirma el intercambio. El retiro por consumo y el cobro completan ese recorrido sin mezclar energía, tokens, financiación y pagos.

Esa conclusión se limita a la evidencia disponible y reconoce las decisiones manuales y autoridades de confianza. El valor del proceso consiste tanto en la integración lograda como en haber delimitado qué no puede inferirse de ella.

\appendix
\apendice{Funciones principales de los contratos}
Este anexo describe las funciones de negocio y grupos heredados relevantes. Los getters generados por mappings y variables públicas, constructores y detalles internos se consultan en el código Solidity y ABI de la revisión fijada. No se presenta esta síntesis como reemplazo de una especificación formal del contrato.

\subsection{TimelockController}
\begin{longtable}{L{.37\linewidth}L{.53\linewidth}}
\toprule Función o grupo & Uso \\ \midrule\endfirsthead
\toprule Función o grupo & Uso \\ \midrule\endhead
\code{schedule}, \code{scheduleBatch} & Programar una acción o grupo con predecesor, salt y retraso.\\
\code{execute}, \code{executeBatch} & Ejecutar lo programado cuando está listo y los roles lo permiten.\\
\code{cancel} & Cancelar operaciones pendientes con la autoridad correspondiente.\\
\code{hashOperation}, \code{hashOperationBatch} & Calcular identificadores deterministas de operaciones.\\
\code{getMinDelay}, \code{updateDelay} & Consultar el retraso o modificarlo mediante la propia gobernanza.\\
Consultas y roles & Comprobar si existe, está pendiente, lista o terminada una operación; consultar, otorgar, revocar o renunciar a roles.\\
\bottomrule
\end{longtable}
\subsection{CommunityRegistry}
\begin{longtable}{L{.37\linewidth}L{.53\linewidth}}
\toprule Función o grupo & Uso \\ \midrule\endfirsthead
\toprule Función o grupo & Uso \\ \midrule\endhead
\code{propose}, \code{transition} & Crear propuesta y aplicar transiciones autorizadas de comunidad y operador.\\
\code{configurePrice}, \code{bindProtocol} & Fijar precio bajo restricciones de estado y vincular contratos válidos.\\
\code{applyMembership}, \code{setMembership} & Solicitar y administrar pertenencia de una wallet.\\
\code{registerNode}, \code{authorizeMeter} & Asociar nodo a propietario y autorizar o revocar medidor.\\
\code{isActive}, \code{isMember} & Consultar elegibilidad de comunidad y participante.\\
\code{pause}, \code{unpause} & Aplicar o retirar pausa con roles correspondientes; AccessControl añade administración de permisos.\\
\bottomrule
\end{longtable}
\subsection{SUNToken}
\begin{longtable}{L{.37\linewidth}L{.53\linewidth}}
\toprule Función o grupo & Uso \\ \midrule\endfirsthead
\toprule Función o grupo & Uso \\ \midrule\endhead
\code{mint} & Emitir con MINTER\_ROLE, cantidad positiva y crédito único.\\
\code{anchorAccounting} & Vincular una referencia de evidencia a un crédito emitido, una sola vez.\\
\code{balanceOf}, \code{totalSupply}, \code{transfer} & Consultar saldo/oferta y transferir unidades existentes.\\
\code{approve}, \code{allowance}, \code{transferFrom} & Autorizar gasto, consultar su límite y usar autorización.\\
\code{burn}, \code{burnFrom} & Quemar unidades propias o autorizadas; una quema genérica no equivale a compensación protocolaria.\\
\code{pause}, \code{unpause}, roles & Controlar pausa, permisos y soporte de interfaces; \code{_update} aplica reglas heredadas a movimientos.\\
\bottomrule
\end{longtable}
\subsection{CommunitySettlement}
\begin{longtable}{L{.37\linewidth}L{.53\linewidth}}
\toprule Función o grupo & Uso \\ \midrule\endfirsthead
\toprule Función o grupo & Uso \\ \midrule\endhead
\code{submit} & Presentar un lote con compromiso y cantidades por la autoridad del oráculo.\\
\code{finalize}, \code{reject} & Resolver un lote con la autoridad y periodo de revisión requeridos.\\
\code{supersede} & Reemplazar un lote no utilizado bajo las restricciones que impiden corregir derechos ya consumidos.\\
\code{issue} & Verificar prueba y autorización y emitir sin superar el total aprobado.\\
\code{retire} & Verificar consumo entregado completo, evitar repeticiones y quemar SUN con allowance.\\
\code{leaf}, \code{bindVault}, \code{allocationRoot} & Calcular hoja, vincular una sola custodia válida y exponer compromiso finalizado.\\
Validaciones internas & \code{_check} y \code{_financiallyUsed} comprueban autorización y uso económico previo.\\
\bottomrule
\end{longtable}
\subsection{NativeSettlementVault}
\begin{longtable}{L{.37\linewidth}L{.53\linewidth}}
\toprule Función o grupo & Uso \\ \midrule\endfirsthead
\toprule Función o grupo & Uso \\ \midrule\endhead
\code{deposit} & Depositar ETH nativo con precio máximo autorizado.\\
\code{depositSUN} & Transferir SUN al escrow con aprobación previa.\\
\code{reserve} & Reservar ambos activos para una asignación probada; exige autoridad de liquidación.\\
\code{execute} & Entregar SUN y acreditar ingreso atómicamente, si la reserva sigue habilitada; puede invocarla cualquier cuenta.\\
\code{withdraw}, \code{withdrawSUN} & Retirar ETH disponible o SUN no reservados.\\
\code{withdrawEarnings} & Cobrar ETH ganado; no convierte SUN ni genera otro derecho.\\
\code{cancel} & Cancelar reserva por el consumidor bajo suspensión o vencimiento del plazo de siete días.\\
\code{leaf} y consultas & Construir hoja y consultar disponibles, reservas, límites y energía entregada por nodo. La función interna \code{_pay} comprueba el envío ETH.\\
\bottomrule
\end{longtable}

\apendice{Consultas de seguimiento y verificación}
Las consultas siguientes corresponden al esquema inspeccionado y no alteran registros. Deben ejecutarse contra la base correcta, sin imprimir cadenas de conexión ni secretos.

\subsection{Recepción y calidad}
\begin{lstlisting}[language=SQL]
SELECT now() AS hora_bd,
       max(measured_at) AS ultima_medicion,
       max(received_at) AS ultima_recepcion,
       count(*) FILTER (
         WHERE received_at > now() - interval '5 minutes'
       ) AS recibidos_5_min
FROM telemetry_messages;

SELECT d.device_code, t.quality_status,
       count(*) AS mensajes, max(t.measured_at) AS ultima
FROM telemetry_messages t
JOIN devices d ON d.id = t.device_id
WHERE t.received_at > now() - interval '10 minutes'
GROUP BY d.device_code, t.quality_status
ORDER BY d.device_code, t.quality_status;
\end{lstlisting}

\subsection{Intervalos elegibles posteriores a recuperación}
\begin{lstlisting}[language=SQL]
SELECT e.id, e.interval_start, e.eligible_export_wh
FROM energy_intervals e
JOIN community_nodes n ON n.home_id = e.home_id
JOIN community_native_deployments d
  ON d.community_id = n.community_id
WHERE e.quality_status = 'FINALIZED'
  AND e.eligible_export_wh > 0
  AND (d.accounting_start IS NULL
       OR e.interval_start >= d.accounting_start)
ORDER BY e.interval_start DESC;
\end{lstlisting}
Este resultado identifica candidatos de evidencia; aún deben pasar política, lote y finalización. No es una consulta de tokens ya emitidos.

\subsection{Estado de lotes, claims y asignaciones}
\begin{lstlisting}[language=SQL]
SELECT id, interval_start, status, archived
FROM native_settlement_batches
ORDER BY interval_start DESC LIMIT 20;

SELECT b.id AS lote, b.status AS estado_lote, b.archived,
       c.kind, c.energy_wh, c.status AS estado_claim
FROM native_settlement_claims c
JOIN native_settlement_batches b ON b.id = c.batch_id
ORDER BY b.interval_start DESC LIMIT 30;

SELECT a.id, a.energy_wh, a.due_wei,
       a.status, a.funding_status, b.archived
FROM native_interval_allocations a
JOIN native_settlement_batches b ON b.id = a.batch_id
ORDER BY b.interval_start DESC LIMIT 30;
\end{lstlisting}
\code corresponde a emisión y \code a retiro por consumo. Los datos SQL son proyecciones y se contrastan con recibos/eventos del despliegue activo antes de afirmar éxito económico.

\apendice{Reproducción del informe y fuentes de evidencia}
\subsection{Compilación}
El proyecto LaTeX usa XeLaTeX, latexmk, Polyglossia, TikZ, fontspec, natbib, booktabs, longtable, caption y listings. No necesita shell-escape, descarga de imágenes ni una base de datos. El cuerpo se configura a 12 puntos con Times New Roman cuando está disponible o Nimbus Roman, una tipografía compatible con Times, como alternativa. Se utilizan márgenes de 2,54 cm, interlineado doble, sangría inicial de 1,27 cm y citas de autor y año. Las tablas utilizan interlineado sencillo para favorecer su lectura.
\begin{lstlisting}[language=bash]
cd docs/informe
make pdf
# Resultado: build/informe.pdf
\end{lstlisting}
Como alternativa, importar la carpeta completa a Overleaf, seleccionar XeLaTeX y definir \code{informe.tex} como documento principal. La portada se personaliza en \code{datos_academicos.tex}: incluye tres autores, institución, la asignatura Tecnologías emergentes, docente y fecha. La fecha y revisión del informe son explícitas; no cambian al recompilar. El PDF es un resultado generado y los archivos auxiliares se excluyen de Git.

\subsection{Qué se conserva en el paquete}
\begin{itemize}
\item Fuentes LaTeX, apartados, referencias y Makefile de compilación.
\item Manifiesto de revisión, tablas, cabecera Alembic, OpenAPI y campañas de prueba, sin secretos.
\item Copia del escenario local público versionado, con procedencia y hash del archivo original.
\item Resumen de la suite seleccionada ejecutada para la revisión documental.
\end{itemize}
El manifiesto \code{evidencia/manifest.json} distingue evidencia versionada de registros locales de trabajo que no forman parte del repositorio general. Los resúmenes de campañas anteriores no convierten esos logs en un conjunto público completo.

\subsection{Glosario}
\begin{longtable}{L{.24\linewidth}L{.66\linewidth}}
\toprule Término & Significado en este proyecto \\ \midrule\endfirsthead
\toprule Término & Significado en este proyecto \\ \midrule\endhead
DePIN & Capa de infraestructura y nodos de medición; en la validación las fuentes son simuladas.\\
DeFi & Reglas de activos, custodia y liquidación mediante contratos; no implica descentralización total.\\
EVM / Anvil & Entorno de ejecución de contratos y nodo local de prueba.\\
ERC-20 / SUN & Estándar de token fungible y crédito energético del protocolo.\\
ETH / wei & Activo nativo y unidad mínima; un ETH equivale a $10^{18}$ wei.\\
Claim & Derecho individual comprometido y probado para emisión o retiro; no es un pago por sí mismo.\\
Escrow / vault & Custodia de activos y contrato que administra disponibilidad, reservas y salidas.\\
Nonce de wallet & Número de secuencia de transacción; distinto del nonce del reto de propiedad.\\
Raíz Merkle & Compromiso resumido del conjunto de hojas; la prueba comprueba pertenencia, no verdad física.\\
Génesis & Identidad del bloque inicial que diferencia cadenas locales con el mismo chain ID.\\
Outbox & Trabajo durable vinculado a una transacción de negocio para evitar pérdida y duplicación.\\
Conciliación & Contraste de proyecciones con recibos/eventos canónicos; no restaura una cadena borrada.\\
\bottomrule
\end{longtable}

% Referencias APA 7, autor y ano.
\begin{thebibliography}{99}
\bibitem[0x3ef(2026a)]{architecture} 0x3ef. (2026a). \textit{Architecture} [Arquitectura]. GitHub. \url{https://github.com/0x3ef/Decentralized-Community-Energy-Network/blob/master/docs/architecture.md}
\bibitem[0x3ef(2026b)]{allocation} 0x3ef. (2026b). \textit{Asignación proporcional mediante mayores residuos} [Código fuente]. GitHub. \url{https://github.com/0x3ef/Decentralized-Community-Energy-Network/blob/master/apps/backend/src/modules/community/allocation.py}
\bibitem[0x3ef(2026c)]{blockchain} 0x3ef. (2026c). \textit{Blockchain protocol} [Protocolo blockchain]. GitHub. \url{https://github.com/0x3ef/Decentralized-Community-Energy-Network/blob/master/docs/blockchain.md}
\bibitem[0x3ef(2026d)]{settlement} 0x3ef. (2026d). \textit{Community energy allocation and settlement} [Asignación y liquidación energética comunitaria]. GitHub. \url{https://github.com/0x3ef/Decentralized-Community-Energy-Network/blob/master/docs/community-settlement.md}
\bibitem[0x3ef(2026e)]{config} 0x3ef. (2026e). \textit{Configuration} [Configuración]. GitHub. \url{https://github.com/0x3ef/Decentralized-Community-Energy-Network/blob/master/docs/configuration.md}
\bibitem[0x3ef(2026f)]{contracts} 0x3ef. (2026f). \textit{Contratos del protocolo energético nativo} [Código fuente]. GitHub. \url{https://github.com/0x3ef/Decentralized-Community-Energy-Network/tree/master/contracts/src/}
\bibitem[0x3ef(2026g)]{database} 0x3ef. (2026g). \textit{Database} [Base de datos]. GitHub. \url{https://github.com/0x3ef/Decentralized-Community-Energy-Network/blob/master/docs/database.md}
\bibitem[0x3ef(2026h)]{software} 0x3ef. (2026h). \textit{Decentralized-Community-Energy-Network} [Código fuente]. GitHub. \url{https://github.com/0x3ef/Decentralized-Community-Energy-Network/tree/master}
\bibitem[0x3ef(2026i)]{deployment} 0x3ef. (2026i). \textit{Deployment and operations} [Despliegue y operaciones]. GitHub. \url{https://github.com/0x3ef/Decentralized-Community-Energy-Network/blob/master/docs/deployment.md}
\bibitem[0x3ef(2026j)]{readme} 0x3ef. (2026j). \textit{Descripción, alcance y limitaciones del proyecto} [README]. GitHub. \url{https://github.com/0x3ef/Decentralized-Community-Energy-Network/blob/master/README.md}
\bibitem[0x3ef(2026k)]{development} 0x3ef. (2026k). \textit{Development and tests} [Desarrollo y pruebas]. GitHub. \url{https://github.com/0x3ef/Decentralized-Community-Energy-Network/blob/master/docs/development.md}
\bibitem[0x3ef(2026l)]{start} 0x3ef. (2026l). \textit{Getting started} [Guía de inicio]. GitHub. \url{https://github.com/0x3ef/Decentralized-Community-Energy-Network/blob/master/docs/getting-started.md}
\bibitem[0x3ef(2026m)]{api} 0x3ef. (2026m). \textit{HTTP API} [API HTTP]. GitHub. \url{https://github.com/0x3ef/Decentralized-Community-Energy-Network/blob/master/docs/api.md}
\bibitem[0x3ef(2026n)]{iot} 0x3ef. (2026n). \textit{IoT metering, simulation and verification} [Medición IoT, simulación y verificación]. GitHub. \url{https://github.com/0x3ef/Decentralized-Community-Energy-Network/blob/master/docs/iot.md}
\bibitem[0x3ef(2026o)]{scenario} 0x3ef. (2026o). \textit{Local native scenario} [Escenario nativo local] [Conjunto de datos de validación local]. GitHub. \url{https://github.com/0x3ef/Decentralized-Community-Energy-Network/blob/master/docs/evidence/local-native-scenario.json}
\bibitem[0x3ef(2026p)]{models} 0x3ef. (2026p). \textit{Registro de modelos SQLModel} [Código fuente]. GitHub. \url{https://github.com/0x3ef/Decentralized-Community-Energy-Network/blob/master/apps/backend/src/core/models.py}
\bibitem[0x3ef(2026q)]{security} 0x3ef. (2026q). \textit{Security and trust model} [Seguridad y modelo de confianza]. GitHub. \url{https://github.com/0x3ef/Decentralized-Community-Energy-Network/blob/master/SECURITY.md}
\bibitem[0x3ef(2026r)]{verification} 0x3ef. (2026r). \textit{Verification record} [Registro de verificación]. GitHub. \url{https://github.com/0x3ef/Decentralized-Community-Energy-Network/blob/master/docs/verification.md}
\end{thebibliography}

\end{document}
